\clearpage
\section{墙面之间的多次漫反射}
把一张白纸放在红色地毯旁，日光照到地毯后，一部分红光会落在纸上。纸若反射了其中的一半，旁边的灰墙也有机会收到这部分光。渲染时，探针可以帮我们记住房间里从各个方向到来的光。上一轮先记下地毯送来的红光；下一轮探针射线遇到白纸，就能查出它除了受到日光照射，还从地毯得到多少光。编辑器的「实时效果 → DDGI → 多次漫反射」控制的就是这次查询。在组合光照 Lit 中，把间接光选为 DDGI，也能找到它。

\pair{profiling/probe-feedback/sponza-single.png}{只让墙面反射一次}
     {profiling/probe-feedback/sponza-multi.png}{墙面也反射上一轮探针送来的光}
两图来自 Redflag 的同一条 Sponza 走廊，展示间接光单项，曝光、机位和八条采样射线的方向序列相同。可沿画面中线看前方的地面。在从左上角数的 $(418,417)$ 处，红通道的线性亮度从 0.02347 变为 0.02691。显示图会经过色调映射，这两个数取自之前的光照输出。画面较暗是因为这里单独显示墙面反射，不含点光源直射。

散射到表面的光要乘一次余弦并在半球上积分，这个量是照度 $E$。理想漫反射材质再乘反射率 $\rho$、除以 $\pi$，得到这块表面朝外送出的辐亮度。Redflag 用余弦加权的半球方向发射探针射线，因此它保存的样本平均值相当于 $H=E/\pi$。计算射线命中点 $y$ 返回的光时，设上一轮从附近探针插值得到 $H_{k-1}(y)$，则
\[
   L_k(y) = L_{\mathrm{direct}}(y) + \rho_y\,H_{k-1}(y).
\]
第一轮取 $H_0=0$。例如白纸的红通道反射率是 $0.5$，收到的 $H$ 是 $0.12$，它在原有直接照明之上多送出 $0.06$。这一项已经包含材质反射率，不能再除一次 $\pi$。NVIDIA 的 \texttt{ProbeTraceRGS.hlsl} 在同一位置查探针照度，再乘 $\rho/\pi$；两种写法的区别在于缓存保存的单位。

实时模式每帧算一轮。每个探针的六个方向从同一份旧数组读取结果，将新估计写到另一份。所有探针算完，两份数组再交换。这样左边探针先执行或右边探针先执行，看到的都是上一轮的房间。为减轻低采样的闪烁，从第二轮起，新估计占两成，旧结果保留八成。烘焙固定每个方向 64 条射线，关闭多次反射只算一轮，开启时算四轮，每轮直接采用新估计。改动模式后，已有的烘焙文件需要重新匹配。

\clearpage
\subsection{三块墙的反射量}
画面变亮有许多原因。为了只测同时经过两个反光面的那部分，测试场景用 A 表示红色地面，B 表示浅灰侧墙，镜头对准后墙 C。灯的位置固定。依次将 A 和 B 染黑，各算一次，再算两者同时变黑和同时保留原色的结果。四种情况的几何和遮挡完全相同，改变的只有材质反射率。用 $I_{AB}$ 表示两面都能反光时，C 中心那个像素的间接光；下式扣掉分别只依赖 A 或只依赖 B 的部分。
\[
     \Delta_{AB}= I_{AB} - I_{A0} - I_{0B} + I_{00}.
\]
一次反射只有一处命中点，不能同时依赖 A 与 B，所以这个差应当为零。第二帧开启反馈时，四个红通道的量分别是 $0.211861$、$0.157684$、$0.049218$ 和 $0$。两条路线都经过 A、B 后送到 C 的光合计约 $0.004959$。它包含两种先后次序。表中均为色调映射前的线性红通道。

\begin{center}\begin{tabular}{lrr}
\toprule
同一块后墙的中心 & 单次反射 $\Delta_{AB}$ & 多次反射 $\Delta_{AB}$\\
\midrule
第一帧，探针还没有上一轮的光 & $<2\times10^{-8}$ & $<2\times10^{-8}$\\
第二帧 & $<2\times10^{-8}$ & 0.004959\\
第三帧 & $<2\times10^{-8}$ & 0.009141\\
第四帧 & $<2\times10^{-8}$ & 0.012139\\
保存静态探针 & $<2\times10^{-8}$ & 0.029409\\
\bottomrule
\end{tabular}\end{center}

下面是 \verb|Probes.glsl| 中 \verb|updateProbe| 的伪代码，每一项都沿当前射线命中的那一面墙计算。旧探针中的 $H$ 由 \verb|gatherDiffuseLight| 从邻近探针取出，方向和距离共同决定各探针的权重。
\begin{verbatim}
hit = traceFromProbe(direction)
if hit:
    L = directLight(hit)                   // 该点收到的直射光经墙面反射
    if feedback and round > 0:
        L += hit.albedo * readOldH(hit)    // 墙面把上一轮收到的光继续反射
    sum += L
writeNewProbe(sum / rayCount)              // 未击中的射线也计入总条数
\end{verbatim}

测试中的反射率最高为 0.68。灯增强 32 倍后，第四帧后墙中心的红通道从 0.228865 变为 7.32368，途中移动再复原相机仍得到这个值。实现将反馈的材质反射率限制在物理范围 $[0,1]$，保留大于 1 的 HDR 光照。NVIDIA 的示例另外把反馈反射率封顶在 0.9；本组有开口的三面墙没有用到这个限制。完全封闭、所有墙都无损反射而且持续发光的房间没有有限的稳态解。

两种模式保持 $12\times6\times7$ 个探针，每个方向 64 条射线。第四次更新发射 193,536 条射线，RTX 5070 Ti Laptop GPU 上，11 次回放的更新 GPU 时间中位数从 0.518 ms 增到 0.683 ms。增加的工作发生在射线命中后，最终从探针给画面着色两次均为 0.0068 ms。数值适用于这个小房间及这组采样数。

{\small \href{https://jcgt.org/published/0008/02/01/}{Majercik 等，JCGT 8(2)，2019，第 1–30 页}\quad
\href{https://github.com/NVIDIAGameWorks/RTXGI-DDGI/blob/f33e496ca31b3f0eec1c4e2cbaa8bb620e337fa6/samples/test-harness/shaders/ddgi/ProbeTraceRGS.hlsl\#L146-L180}{NVIDIA 命中点着色源码}\quad
\href{run:../profiling/probe-feedback/README.md}{逐像素记录、项目笔记与复现命令}}
